\subsection{A nonregular family with locally finite Steenrod operations}

\begin{theorem}\label{E:nilpotent-inflations}
Let $N$ be a finite nilpotent semigroup with zero $0$, and let $I,J$
be finite nonempty sets with chosen elements $i_0,j_0$.  Define
\[
 M=\{1\}\amalg(I\times J)\amalg(N\setminus\{0\})
\]
as follows.  The element $1$ is the identity, the set $I\times J$ is a
rectangular-band ideal, and, writing $e=(i_0,j_0)$,
\[
 (i,j)(i',j')=(i,j'),\qquad
 n(i,j)=(i_0,j),\qquad (i,j)n=(i,j_0).
\]
For $n,n'\in N\setminus\{0\}$, their product is their product in $N$
if that is nonzero, and is $e$ otherwise.
Then $M$ is an aperiodic finite monoid, and for every prime $p$ its
cohomology $H^*(BM;\mathbb F_p)$ is locally finite as a module over the
Steenrod algebra.  In particular $BM$ is not homotopy equivalent to
$\Sigma\CP^\infty$.

More precisely, let $Q=N^1$ be the pointed monoid with zero as basepoint
and a freshly adjoined identity, and let
$\epsilon:Q\to S^0=\{*,1\}$ send all nonunits to the basepoint.
Let $X_N$ be the based two-sided simplicial bar
$B(S^0,Q,S^0)$ for this augmentation, with smash products in its levels.
There is a stable equivalence
\[
 \widetilde{\Sigma^\infty BM}\ \simeq
 \bigvee^{(|I|-1)(|J|-1)}\Sigma^2\Sigma^\infty X_N.       \tag{E.20}
\]
Every reduced cohomology class of $X_N$ is pulled back from a finite
CW quotient of $X_N$.
\end{theorem}

\begin{proof}
\emph{The multiplication and the spherical ideal.}
Associativity of three outside factors follows from associativity in
$N$, with its zero replaced by $e$.  The element $e$ absorbs every
outside factor.  With one or two band factors, the displayed row and
column rules give the same outer coordinates in either association.
These checks cover all triples.  An outside element has a power equal
to $e$, and every band element is idempotent, so all eventual cyclic
groups are trivial.  Thus $M$ is aperiodic.

We give the stable argument explicitly because an assertion only about
the ordinary integral monoid algebra would not control Steenrod
operations.  Write $\mathbb S$ for the sphere spectrum and put
\[
 R=\Sigma^\infty M_+,\qquad
 K=\Sigma^\infty(I\times J)_+,\qquad
 A=\Sigma^\infty Q.
\]
The multiplication of a discrete monoid makes its suspension spectrum
an associative ring spectrum.  The ideal inclusion and quotient give
a cofiber sequence of $R$-bimodules
\[
 K\longrightarrow R\longrightarrow A.                 \tag{E.21}
\]
Here the quotient identifies the entire band with the basepoint; it
does not identify the outside elements with one another.

The left module $Re$ and right module $eR$ are literal retracts of
the regular modules, using multiplication by the idempotent $e$.
Their underlying spectra are the wedges on $I$ and $J$, respectively,
and $eRe=\mathbb S$.  Multiplication induces an equivalence
\[
 Re\mathbin{\wedge}_{\mathbb S}eR\ \xrightarrow{\ \cong\ }K
                                                        \tag{E.22}
\]
of bimodules: on the indexing sets it is the bijection
$((i,j_0),(i_0,j))\mapsto(i,j)$.
Derived tensor products with these representable retracts are their
ordinary idempotent summands.  In particular $Ae=eA=0$.  Tensoring
(E.21) with $A$ and using (E.22) proves
\[
 A\mathbin{\wedge}^{\mathbb L}_R A\simeq A.             \tag{E.23}
\]
Consequently restriction from $A$-modules to $R$-modules is fully
faithful: its extension--restriction counit is an equivalence by
(E.23) and associativity of derived tensor products.

\emph{The two augmentation fibers.}
Regard $\mathbb S$ as a trivial left or right $R$-module, so every
element of $M$ acts as the identity.  Let
\[
 V=\operatorname{fib}(Re\longrightarrow\mathbb S),\qquad
 W=\operatorname{fib}(eR\longrightarrow\mathbb S),        \tag{E.24}
\]
where the maps fold the indicated wedges.  Applying $e$ to the first
fiber gives the fiber of the literal identity
$\mathbb S\to\mathbb S$, so $eV\simeq0$; similarly $We\simeq0$.
This is a vanishing of spectra with their module structure, not just
of their degree-zero homotopy groups.  Equation (E.22) therefore gives
$K\wedge_R^{\mathbb L}V\simeq0$.  Tensoring (E.21) with $V$ gives
\[
 V\simeq A\mathbin{\wedge}^{\mathbb L}_R V.
\]
It supplies a coherent $A$-module structure on $V$.  The analogous
argument gives one on $W$.

The usual simplicial bar is a model for derived tensor products, and
in particular
$\mathbb S\wedge_R^{\mathbb L}\mathbb S\simeq\Sigma^\infty_+BM$.
Tensor the first fiber sequence (E.24) with the trivial right module.
Its middle term is $\mathbb S e=\mathbb S$, and the resulting map to
$\Sigma^\infty_+BM$ is split by the augmentation to $\mathbb S$.
Taking the reduced summand gives
\[
 \widetilde{\Sigma^\infty BM}
       \simeq\Sigma(\mathbb S\wedge_R^{\mathbb L}V).
\]
Now tensor the second sequence (E.24) with $V$.  Its middle term is
$eV\simeq0$, whence
\[
 \widetilde{\Sigma^\infty BM}
       \simeq\Sigma^2(W\wedge_R^{\mathbb L}V)
       \simeq\Sigma^2(W\wedge_A^{\mathbb L}V).           \tag{E.25}
\]
The last equivalence follows from (E.23) since both modules descend
to $A$.

\emph{Identification of the modules, including coherence.}
Let $T=\{1,e\}\amalg(N\setminus\{0\})\subset M$.
The element $e$ is a zero in $T$.  Put $P=\Sigma^\infty T_+$.
The same one-element ideal calculation gives
$A\wedge_P^{\mathbb L}A\simeq A$, so restriction of $A$-modules to
$P$-modules is fully faithful as well.

On restriction to $T$, the row $(i_0,j_0)$ defines a module section
$\mathbb S\to Re$ of its augmentation: it is fixed by every element
of $T$.  In a stable category the fiber of a split epimorphism is the
cofiber of its section.  Here that cofiber is the strict pointed
permutation quotient $\Sigma^\infty(I/\{i_0\})$.
Every nonunit of $T$ sends all rows to $i_0$, hence acts as the zero
map on this quotient.  Thus, as a $P$-module, it is precisely a wedge
of $|I|-1$ copies of the sphere module defined by $\epsilon:A\to\mathbb S$.
Both it and $V$ are restrictions of $A$-modules.  Full faithfulness
therefore lifts their equivalence to an equivalence of $A$-modules.
The column argument gives, with the proper right action,
\[
 V\simeq\bigvee^{|I|-1}\mathbb S_\epsilon,
 \qquad W\simeq\bigvee^{|J|-1}\mathbb S_\epsilon.
\]
This argument uses actual module maps and the two quotient equivalences;
no unproved lift of a degree-zero idempotent or matrix is used.
The based simplicial bar identifies
$\mathbb S_\epsilon\wedge_A^{\mathbb L}\mathbb S_\epsilon$
with $\Sigma^\infty X_N$.  Substitution in (E.25) proves (E.20).

\emph{Finite quotients detecting every cohomology class.}
Choose $\ell\geq2$ with $N^\ell=\{0\}$.  For $a\ne0$ set
\[
 w(a)=\max\{r:a\in N^r\}.
\]
Then $1\leq w(a)\leq\ell-1$, and if $ab\ne0$ then
$w(ab)\geq w(a)+w(b)$.  A nonbasepoint nondegenerate positive-dimensional
simplex of $X_N$ is a word $(a_1,\ldots,a_n)$ of nonzero elements of $N$.
Its two exterior faces are the basepoint.  An interior face multiplies
two adjacent letters, and is the basepoint if their product is zero.
Otherwise it does not decrease the total weight
$\sum_iw(a_i)$.  Degeneracies insert the identity, of weight zero.

For an integer $q\geq1$, the basepoint and all words of total weight
at least $q$ consequently form a based simplicial subspace $F_qX_N$.
Every nondegenerate cell of this subspace has dimension at least
$\lceil q/(\ell-1)\rceil$.  In contrast, every nondegenerate positive
cell of the quotient $X_N/F_qX_N$ has dimension at most $q-1$,
because every letter has weight at least one.  The quotient is a finite
CW complex, since $N$ is finite.  Its nonbasepoint zero-simplex also
causes no problem and is retained in the quotient.

Given a fixed degree $d$, choose $q>(d+1)(\ell-1)$.
The subspace $F_qX_N$ has no positive cells in dimensions at most
$d+1$.  The cofiber sequence shows that
\[
 \widetilde H^d(X_N/F_qX_N;\mathbb F_p)
     \longrightarrow\widetilde H^d(X_N;\mathbb F_p)
\]
is an isomorphism.  Thus every class lifts to a finite quotient.
Naturality of Steenrod operations makes the entire Steenrod orbit
of that class an image of the finite-dimensional cohomology of this
one quotient.  It is therefore finite-dimensional.  The same holds
for (E.20), since Steenrod operations commute with suspension and
stable equivalences.

Finally, if $u\in H^2(\CP^\infty;\mathbb F_p)$ is the generator,
its suspension has nonzero successive operations
$P^{p^{r-1}}\cdots P^1(\sigma u)=\sigma(u^{p^r})$
at every odd prime.  At $p=2$ the explicit sequence is
$Sq^{2^r}\cdots Sq^4 Sq^2(\sigma u)=\sigma(u^{2^r})$ for $r\geq1$,
where the first operation is $Sq^2$ and the intermediate string is
empty when $r=1$.  Its Steenrod orbit is infinite-dimensional,
which proves the asserted exclusion.
\end{proof}

\begin{remark}\label{E:nilpotent-inflations-scope}
The theorem includes the inflated rectangular-band infinite-wedge
example (one nonzero nilpotent letter with square zero).  Its proof
does not assume bounded homology or finite-dimensional homotopy type.
It proves local finiteness for the stated family, not for all aperiodic
monoids or all finite categories.  The next subsection removes the
constant-action hypothesis and allows a finite unit group at suitable
primes.  Further nonzero idempotents in a general Rees quotient
require a separate argument; the controlled-image localization and
multibar induction of Theorem~\ref{E:general-controlled-multibar}
provide it below.
\end{remark}
